% Figure peps_patch_first_variation.
% Owner: chapter/ch24_peps_regularized_patch_stationarity.tex (thm:peps_patch_first_variation).
% Used by: gapped-2d-notes PP Sec. 5 step (b); relabelled there as
% L_{>j} [B,L_j] L_{<j} Omega, the derivative M'(0)Omega (not the paper's
% log-derivative insertion B - L_j B L_j^{-1}).
% Formula: (Delta_j)_sigma = sum_{alpha,beta,gamma}
% (P_j^+)_{sigma,alpha} (D_j)_{alpha,beta}
% (P_j^-)_{beta,gamma} Omega_gamma.
% Source: 03-patches.tex lines 134--149, adc7f1241b42e322a6451854ab7e4b4c146bf78a.
% Ink-to-index map: every wire is a global configuration index.
% D_j is the identity extension of the regional commutator; P_j^+ and
% P_j^- are global ordered products, not operators on X_j.
% Contractions: alpha, beta, gamma, in that left-to-right order.
% Open leg: sigma on the left. Boundary signature: phys:w.
% Action begins at Omega on the right: P_j^-, then D_j, then P_j^+.
% No tensor-factor decomposition, nesting, or commutation is assumed.
% Label macros: \tnlFvOuter{P_j^+}, \tnlFvInsert{D_j}, \tnlFvInner{P_j^-};
% the PEPS note overrides them with L_{>j}, [B,L_j], L_{<j}.
\begin{center}
    \tenkzkernel
    \begin{tenkz}[rows={wire}, cols=7, bonds=none]
        \providecommand{\tnlFvOuter}{P_j^+}
        \providecommand{\tnlFvInsert}{D_j}
        \providecommand{\tnlFvInner}{P_j^-}
        \tn[at=(1,1), name=coordinateouter, skin=box,
            ports={180:physical, 0:physical}]{\tnlFvOuter}
        \tn[at=(1,3), name=coordinateinsert, skin=box,
            ports={180:physical, 0:physical}]{\tnlFvInsert}
        \tn[at=(1,5), name=coordinateinner, skin=box,
            ports={180:physical, 0:physical}]{\tnlFvInner}
        \tn[at=(1,7), name=coordinatestate, skin=box,
            ports={180:physical}]{\Omega}
        \tnwire[name=coordinateoutput]{open w}{coordinateouter.180}
        \tnwire[name=coordinatealpha]{coordinateouter.0}{coordinateinsert.180}
        \tnwire[name=coordinatebeta]{coordinateinsert.0}{coordinateinner.180}
        \tnwire[name=coordinategamma]{coordinateinner.0}{coordinatestate.180}
        \tnmark[form=label, label pos=n]{on coordinateoutput 0.5}{$\sigma$}
        \tnmark[form=label, label pos=n]{on coordinatealpha 0.5}{$\alpha$}
        \tnmark[form=label, label pos=n]{on coordinatebeta 0.5}{$\beta$}
        \tnmark[form=label, label pos=n]{on coordinategamma 0.5}{$\gamma$}
    \end{tenkz}
\end{center}
